\documentclass[11pt,a4paper]{article}
\usepackage[utf8]{inputenc}
\usepackage[margin=1in]{geometry}
\usepackage{amsmath,amssymb,amsfonts}
\usepackage{booktabs}
\usepackage{graphicx}
\usepackage{xcolor}
\usepackage{listings}
\usepackage{hyperref}
\usepackage{float}
\usepackage{longtable}
\usepackage{array}
\usepackage{adjustbox}

% Code listing style
\definecolor{codegreen}{rgb}{0,0.6,0}
\definecolor{codegray}{rgb}{0.5,0.5,0.5}
\definecolor{codepurple}{rgb}{0.58,0,0.82}
\definecolor{backcolour}{rgb}{0.95,0.95,0.92}

\lstdefinestyle{pythonstyle}{
    backgroundcolor=\color{backcolour},
    commentstyle=\color{codegreen},
    keywordstyle=\color{blue},
    numberstyle=\tiny\color{codegray},
    stringstyle=\color{codepurple},
    basicstyle=\ttfamily\footnotesize,
    breakatwhitespace=false,
    breaklines=true,
    captionpos=b,
    keepspaces=true,
    numbers=left,
    numbersep=5pt,
    showspaces=false,
    showstringspaces=false,
    showtabs=false,
    tabsize=2,
    language=Python
}

\lstset{style=pythonstyle}

\title{\textbf{QPU Benchmark Modes: \\
Detailed Analysis and Code-to-Equation Verification}\\[0.5em]
\large Technical Report for \texttt{run\_qpu\_benchmarks\_all\_scenarios.py}}
\author{OQI-UC002-DWave Project}
\date{February 2026}

\begin{document}

\maketitle

\tableofcontents
\newpage

%%%%%%%%%%%%%%%%%%%%%%%%%%%%%%%%%%%%%%%%%%%%%%%%%%%%%%%%%%%%%%%%%%%%%%%%%%%%%%%
\section{Executive Summary}
%%%%%%%%%%%%%%%%%%%%%%%%%%%%%%%%%%%%%%%%%%%%%%%%%%%%%%%%%%%%%%%%%%%%%%%%%%%%%%%

This report provides a comprehensive analysis of the three Quantum Processing Unit (QPU) modes implemented in the unified benchmark system for crop rotation optimization. All modes implement \textbf{Variant B: Hierarchical Multi-Period Rotation} as defined in the main project report.

\begin{enumerate}
    \item \textbf{qpu-native-6-family}: Direct 6-family BQM without aggregation
    \item \textbf{qpu-hierarchical-aggregated}: 27$\to$6 aggregation with spatial decomposition and refinement
    \item \textbf{qpu-hybrid-27-food}: Full 27-food variables with 6-family synergy template
\end{enumerate}

Each mode represents a different trade-off between problem fidelity and QPU compatibility.

%%%%%%%%%%%%%%%%%%%%%%%%%%%%%%%%%%%%%%%%%%%%%%%%%%%%%%%%%%%%%%%%%%%%%%%%%%%%%%%
\section{Problem Formulation: Variant B Multi-Period Rotation}
%%%%%%%%%%%%%%%%%%%%%%%%%%%%%%%%%%%%%%%%%%%%%%%%%%%%%%%%%%%%%%%%%%%%%%%%%%%%%%%

All QPU modes solve the \textbf{Variant B} formulation from the main project report. This section reproduces the mathematical formulation for direct comparison.

\subsection{Decision Variables}

Let $Y_{f,c,t}$ be a binary variable:
\begin{equation}
Y_{f,c,t} = \begin{cases} 
1 & \text{if crop } c \in \mathcal{C} \text{ is planted on farm } f \in \mathcal{F} \text{ in period } t \in \mathcal{T} \\
0 & \text{otherwise}
\end{cases}
\end{equation}

\subsection{Objective Function}

The objective is to maximize the total normalized benefit:
\begin{equation}
\max \quad \left( \text{Benefit} + \text{Temporal} + \text{Spatial} - \text{Penalty} + \text{Diversity} \right)
\end{equation}

\subsubsection{Base Benefit (Linear)}
\begin{equation}
\text{Benefit} = \sum_{f \in \mathcal{F}} \sum_{c \in \mathcal{C}} \sum_{t \in \mathcal{T}} \frac{B_c \cdot L_f}{A_{\text{total}}} \cdot Y_{f,c,t}
\end{equation}

where $B_c$ is the composite benefit score for crop $c$, $L_f$ is the land availability for farm $f$, and $A_{\text{total}} = \sum_{f} L_f$.

\subsubsection{Temporal Synergy (Quadratic)}
\begin{equation}
\text{Temporal} = \gamma_{\text{rot}} \sum_{f \in \mathcal{F}} \sum_{t=2}^{T} \sum_{c_1 \in \mathcal{C}} \sum_{c_2 \in \mathcal{C}} \frac{L_f}{A_{\text{total}}} \cdot R_{c_1,c_2} \cdot Y_{f,c_1,t-1} \cdot Y_{f,c_2,t}
\end{equation}

The rotation matrix $R \in \mathbb{R}^{|\mathcal{C}| \times |\mathcal{C}|}$ is defined as:
\begin{equation}
R_{c_1,c_2} = \begin{cases}
-\beta \cdot 1.5 & \text{if } c_1 = c_2 \text{ (monoculture penalty)} \\
\text{Unif}(\beta \cdot 1.2, \beta \cdot 0.3) & \text{with prob. } p_{\text{frust}} \text{ (disease/competition)} \\
\text{Unif}(0.02, 0.20) & \text{otherwise (beneficial rotation)}
\end{cases}
\end{equation}

\subsubsection{Spatial Synergy (Quadratic)}
\begin{equation}
\text{Spatial} = \gamma_{\text{spat}} \sum_{(f_1, f_2) \in \mathcal{E}} \sum_{t \in \mathcal{T}} \sum_{c_1 \in \mathcal{C}} \sum_{c_2 \in \mathcal{C}} \frac{S_{c_1,c_2}}{A_{\text{total}}} \cdot Y_{f_1,c_1,t} \cdot Y_{f_2,c_2,t}
\end{equation}
where $\mathcal{E}$ is the edge set of the $k$-nearest neighbor graph and $S_{c_1,c_2} = 0.3 \cdot R_{c_1,c_2}$.

\subsubsection{One-Hot Penalty (Quadratic)}
\begin{equation}
\text{Penalty} = \lambda_{\text{OH}} \sum_{f \in \mathcal{F}} \sum_{t \in \mathcal{T}} \left( \left(\sum_{c \in \mathcal{C}} Y_{f,c,t}\right) - 1 \right)^2
\end{equation}

\subsubsection{Diversity Bonus (Linear)}
\begin{equation}
\text{Diversity} = \lambda_{\text{div}} \sum_{f \in \mathcal{F}} \sum_{c \in \mathcal{C}} \mathbb{I}\left(\sum_{t \in \mathcal{T}} Y_{f,c,t} > 0\right)
\end{equation}

\subsection{Model Parameters}

\begin{table}[H]
\centering
\caption{Variant B model parameters (from main report Table 8)}
\label{tab:variant_b_params}
\begin{tabular}{@{}lcc@{}}
\toprule
\textbf{Parameter} & \textbf{Symbol} & \textbf{Default Value} \\
\midrule
Negative synergy strength & $\beta$ & $-0.8$ \\
Frustration ratio & $p_{\text{frust}}$ & $0.70$ \\
Temporal rotation weight & $\gamma_{\text{rot}}$ & $0.2$ \\
Spatial coupling strength & $\gamma_{\text{spat}}$ & $0.5 \gamma_{\text{rot}} = 0.1$ \\
Spatial dampening & --- & $0.3$ (applied to $S$) \\
One-hot penalty & $\lambda_{\text{OH}}$ & $3.0$ \\
Diversity bonus & $\lambda_{\text{div}}$ & $0.15$ \\
Planning periods & $T$ & $3$ \\
\bottomrule
\end{tabular}
\end{table}

\subsection{BQM Conversion}

For quantum annealing, the objective is negated (since QA minimizes) and expressed in BQM form:
\begin{equation}
E(\mathbf{Y}) = \sum_{(f,c,t)} h_{f,c,t} \cdot Y_{f,c,t} + \sum_{\substack{(f,c,t) \\ (f',c',t')}} J_{(f,c,t),(f',c',t')} \cdot Y_{f,c,t} \cdot Y_{f',c',t'}
\end{equation}

where $h_{f,c,t}$ are linear biases and $J$ encodes quadratic couplings.

%%%%%%%%%%%%%%%%%%%%%%%%%%%%%%%%%%%%%%%%%%%%%%%%%%%%%%%%%%%%%%%%%%%%%%%%%%%%%%%
\section{QPU Mode Implementations}
%%%%%%%%%%%%%%%%%%%%%%%%%%%%%%%%%%%%%%%%%%%%%%%%%%%%%%%%%%%%%%%%%%%%%%%%%%%%%%%

\subsection{Mode 1: qpu-native-6-family}

\subsubsection{Description}
Solves the 6-family problem directly on QPU without aggregation. For scenarios with 6 food families, this is the most direct approach.

\subsubsection{Linear Bias Construction}

From the BQM conversion of the Variant B objective:
\begin{equation}
h_{f,c,t} = -\frac{B_c \cdot L_f}{A_{\text{total}}} - \frac{\lambda_{\text{div}}}{T} - \lambda_{\text{OH}}
\end{equation}

\begin{lstlisting}[caption={Linear Bias Implementation}]
# Linear bias: negative benefit (BQM minimizes)
linear_bias = -benefit * area_frac
# Add diversity bonus
linear_bias -= diversity_bonus / n_periods
# Add one-hot linear part
linear_bias -= one_hot_penalty

bqm.add_variable(var_name, linear_bias)
\end{lstlisting}

\subsubsection{Temporal Synergy Coupling}

Implements the Temporal term from Variant B:
\begin{equation}
J^{\text{temp}}_{(f,c_1,t-1),(f,c_2,t)} = -\gamma_{\text{rot}} \cdot R_{c_1,c_2} \cdot \frac{L_f}{A_{\text{total}}}
\end{equation}

\begin{lstlisting}[caption={Temporal Synergy Implementation}]
# Temporal synergies
for f_idx, farm in enumerate(farm_names):
    area_frac = land_availability[farm] / total_area
    for t in range(2, n_periods + 1):
        for c1_idx in range(n_foods):
            for c2_idx in range(n_foods):
                synergy = R[c1_idx, c2_idx]
                var1 = var_map[(f_idx, c1_idx, t-1)]
                var2 = var_map[(f_idx, c2_idx, t)]
                bqm.add_quadratic(var1, var2, 
                    -rotation_gamma * synergy * area_frac)
\end{lstlisting}

\subsubsection{Spatial Synergy Coupling}

Implements the Spatial term with dampened matrix $S = 0.3 \cdot R$:
\begin{equation}
J^{\text{spat}}_{(f_1,c_1,t),(f_2,c_2,t)} = -\gamma_{\text{spat}} \cdot S_{c_1,c_2} = -0.5\gamma_{\text{rot}} \cdot 0.3 R_{c_1,c_2}
\end{equation}

\begin{lstlisting}[caption={Spatial Synergy Implementation}]
# Spatial synergies (within adjacent farms)
for f_idx in range(n_farms - 1):
    for t in range(1, n_periods + 1):
        for c1_idx in range(n_foods):
            for c2_idx in range(n_foods):
                synergy = R[c1_idx, c2_idx] * 0.3  # S = 0.3*R
                var1 = var_map[(f_idx, c1_idx, t)]
                var2 = var_map[(f_idx + 1, c2_idx, t)]
                bqm.add_quadratic(var1, var2, 
                    -rotation_gamma * 0.5 * synergy)
\end{lstlisting}

\subsubsection{One-Hot Penalty}

Expanding the quadratic penalty term:
\begin{equation}
\left(\sum_c Y_{f,c,t} - 1\right)^2 = \sum_c Y_{f,c,t} - 2\sum_c Y_{f,c,t} + 2\sum_{c<c'} Y_{f,c,t} Y_{f,c',t} + 1
\end{equation}

The quadratic contribution is:
\begin{equation}
J^{\text{OH}}_{(f,c,t),(f,c',t)} = 2\lambda_{\text{OH}} \quad \text{for } c \neq c'
\end{equation}

\begin{lstlisting}[caption={One-Hot Penalty Implementation}]
# One-hot penalty (quadratic part)
for f_idx in range(n_farms):
    for t in range(1, n_periods + 1):
        vars_this = [var_map[(f_idx, c, t)] for c in range(n_foods)]
        for i in range(len(vars_this)):
            for j in range(i + 1, len(vars_this)):
                bqm.add_quadratic(vars_this[i], vars_this[j], 
                    2 * one_hot_penalty)
\end{lstlisting}

\subsubsection{Rotation Matrix Generation}

\begin{lstlisting}[caption={Family Rotation Matrix (matching Variant B Eq. 10)}]
def create_family_rotation_matrix(seed: int = 42) -> np.ndarray:
    np.random.seed(seed)
    n_families = 6
    R = np.zeros((n_families, n_families))
    
    frustration_ratio = 0.7   # p_frust from Variant B
    negative_strength = -0.8  # beta from Variant B
    
    for i in range(n_families):
        for j in range(n_families):
            if i == j:
                # Same family: monoculture penalty = -beta * 1.5
                R[i, j] = negative_strength * 1.5  # = -1.2
            elif np.random.random() < frustration_ratio:
                # Disease/competition: Unif(beta*1.2, beta*0.3)
                R[i, j] = np.random.uniform(
                    negative_strength * 1.2, 
                    negative_strength * 0.3)
            else:
                # Beneficial rotation: Unif(0.02, 0.20)
                R[i, j] = np.random.uniform(0.02, 0.20)
    return R
\end{lstlisting}

\subsection{Mode 2: qpu-hierarchical-aggregated}

\subsubsection{Three-Level Pipeline}

\begin{enumerate}
    \item \textbf{Level 1 (Classical)}: Aggregate 27 foods $\to$ 6 families
    \item \textbf{Level 2 (Quantum)}: Solve 6-family BQM with spatial decomposition
    \item \textbf{Level 3 (Classical)}: Refine back to 27 foods for MIQP scoring
\end{enumerate}

\subsubsection{Benefit Aggregation}

Benefits are averaged by family:
\begin{equation}
B_g = \frac{1}{|\mathcal{C}_g|} \sum_{c \in \mathcal{C}_g} B_c
\end{equation}

where $\mathcal{C}_g$ is the set of crops in family $g$.

\subsubsection{Spatial Decomposition}

Farms are partitioned into clusters of size $k$:
\begin{equation}
n_{\text{vars}}^{\text{cluster}} = k \times |\mathcal{C}| \times T = k \times 6 \times 3 = 18k
\end{equation}

For $k=9$ (default), each cluster has 162 variables, safely below the QPU limit of $\sim$177.

\subsubsection{Boundary Coordination}

After each iteration, boundary farms receive soft constraints from neighboring clusters:
\begin{equation}
h^{\text{boundary}}_{f,c,t} \leftarrow h_{f,c,t} - 0.5 \gamma_{\text{rot}} \quad \text{if } Y^{\text{prev}}_{f_{\text{neighbor}},c,t} = 1
\end{equation}

\subsection{Mode 3: qpu-hybrid-27-food}

\subsubsection{Hybrid Rotation Matrix}

Uses full 27-food variable space but derives synergy from 6-family template:
\begin{equation}
R^{27}_{c_1,c_2} = R^6_{g(c_1), g(c_2)} + \epsilon_{c_1,c_2}
\end{equation}

where $g(c)$ maps food $c$ to its family and $\epsilon \sim \mathcal{U}(-0.05, 0.05)$.

\subsubsection{Cluster Size Constraints}

With 27 foods $\times$ 3 periods = 81 variables per farm:
\begin{equation}
k_{\max} = \left\lfloor \frac{170}{81} \right\rfloor = 2 \text{ farms per cluster}
\end{equation}

%%%%%%%%%%%%%%%%%%%%%%%%%%%%%%%%%%%%%%%%%%%%%%%%%%%%%%%%%%%%%%%%%%%%%%%%%%%%%%%
\section{Data Sources and Plot Mapping}
%%%%%%%%%%%%%%%%%%%%%%%%%%%%%%%%%%%%%%%%%%%%%%%%%%%%%%%%%%%%%%%%%%%%%%%%%%%%%%%

This section documents which result files are used by \texttt{generate\_single\_plot\_pdfs.py}.

\subsection{Primary Data Files}

\begin{table}[H]
\centering
\caption{Data files used by the plotting script}
\label{tab:data_files}
\begin{adjustbox}{max width=\textwidth}
\begin{tabular}{@{}lll@{}}
\toprule
\textbf{Variable} & \textbf{File Path} & \textbf{Description} \\
\midrule
\texttt{DEFAULT\_QPU\_HIER\_PATH} & \texttt{qpu\_hier\_repaired.json} & QPU hierarchical results (repaired) \\
\texttt{DEFAULT\_GUROBI\_60S\_PATH} & \texttt{gurobi\_baseline\_60s.json} & Gurobi 60s timeout baseline \\
\texttt{DEFAULT\_GUROBI\_300S\_FILE} & \texttt{@todo/.../gurobi\_timeout\_test\_*.json} & Gurobi 300s timeout results \\
\bottomrule
\end{tabular}
\end{adjustbox}
\end{table}

\subsection{Plot Categories and Data Sources}

\begin{table}[H]
\centering
\caption{Mapping of plot sections to data sources}
\label{tab:plot_mapping}
\small
\begin{adjustbox}{max width=\textwidth}
\begin{tabular}{@{}p{4cm}p{3cm}p{6cm}@{}}
\toprule
\textbf{Plot Section} & \textbf{Data Source} & \textbf{Plots Generated} \\
\midrule
Section 1: Scaling & QPU + Gurobi 60s & \texttt{scaling\_objectives\_by\_vars/farms}, \texttt{scaling\_time\_comparison}, \texttt{scaling\_qpu\_breakdown\_by\_vars/farms} \\
\addlinespace
Section 2: Split Formulation & QPU + Gurobi 300s & \texttt{split\_objectives\_by\_vars/farms}, \texttt{split\_optimality\_gap\_by\_vars/farms}, \texttt{split\_solve\_time\_by\_vars/farms}, \texttt{split\_speedup\_by\_vars/farms}, \texttt{split\_pure\_qpu\_time\_by\_vars/farms}, \texttt{split\_gurobi\_mip\_gap\_by\_vars/farms} \\
\addlinespace
Section 3: Gap Analysis & QPU + Gurobi 300s & \texttt{gap\_absolute\_comparison}, \texttt{gap\_ratio}, \texttt{gap\_vs\_mip\_gap}, \texttt{gap\_6family\_scaling}, \texttt{gap\_27food\_scaling}, \texttt{gap\_summary\_table} \\
\addlinespace
Section 4: Violation Impact & Dedicated loader & \texttt{violation\_rate}, \texttt{violation\_gap\_breakdown}, \texttt{violation\_adjusted\_objective} \\
\addlinespace
Section 5: Deep Dive & Dedicated loader & \texttt{deepdive\_objective\_comparison}, \texttt{deepdive\_ratio\_analysis}, \texttt{deepdive\_gap\_attribution}, \texttt{deepdive\_scatter}, \texttt{deepdive\_violation\_impact\_by\_scenario}, \texttt{deepdive\_summary\_findings} \\
\addlinespace
Section 6: QPU Advantage & QPU + Gurobi 300s & \texttt{advantage\_benefit\_comparison}, \texttt{advantage\_benefit\_ratio}, \texttt{advantage\_time\_comparison}, \texttt{advantage\_violations\_vs\_benefit}, \texttt{advantage\_pure\_qpu\_scaling}, \texttt{advantage\_summary\_table} \\
\addlinespace
Section 7: Benchmarks & Benchmark loader & \texttt{benchmark\_small\_*}, \texttt{benchmark\_large\_*}, \texttt{benchmark\_comprehensive\_*} \\
\bottomrule
\end{tabular}
\end{adjustbox}
\end{table}

\subsection{Result Files Analyzed}

Three JSON result files from the \texttt{qpu-hierarchical-aggregated} mode were analyzed:

\begin{enumerate}
    \item \texttt{qpu\_hier\_all\_6family.json}: Hierarchical mode with \texttt{farms\_per\_cluster=5}
    \item \texttt{qpu\_hier\_repaired.json}: Hierarchical mode with repair and \texttt{farms\_per\_cluster=9}
    \item \texttt{qpu\_hierarchical\_all.json}: Hierarchical mode with \texttt{farms\_per\_cluster=10}
\end{enumerate}

%%%%%%%%%%%%%%%%%%%%%%%%%%%%%%%%%%%%%%%%%%%%%%%%%%%%%%%%%%%%%%%%%%%%%%%%%%%%%%%
\section{Results Analysis}
%%%%%%%%%%%%%%%%%%%%%%%%%%%%%%%%%%%%%%%%%%%%%%%%%%%%%%%%%%%%%%%%%%%%%%%%%%%%%%%

\subsection{Objective Quality Comparison}

\begin{table}[H]
\centering
\caption{MIQP Objective Values by Configuration}
\begin{tabular}{llrrr}
\toprule
\textbf{Scenario} & \textbf{Farms} & \textbf{All-6Family} & \textbf{Repaired} & \textbf{All-10} \\
\midrule
rotation\_micro\_25 & 5 & -5.50 & -4.86 & -5.75 \\
rotation\_small\_50 & 10 & -3.88 & -21.79 & error \\
rotation\_15farms\_6foods & 15 & -5.10 & -26.22 & error \\
rotation\_medium\_100 & 20 & -24.89 & -39.24 & error \\
rotation\_25farms\_6foods & 25 & -31.15 & -52.67 & error \\
rotation\_large\_200 & 50 & -37.47 & N/A & error \\
rotation\_50farms\_6foods & 50 & N/A & -109.67 & error \\
\bottomrule
\end{tabular}
\end{table}

\subsection{Key Findings}

\subsubsection{Embedding Errors}

The \texttt{qpu\_hierarchical\_all.json} file shows systematic embedding failures when \texttt{farms\_per\_cluster=10}:
\begin{verbatim}
"error_message": "Cannot embed given BQM (size 180), 
                  sampler can only handle problems of size 177"
\end{verbatim}

This occurs because $10 \times 6 \times 3 = 180$ variables exceeds the D-Wave Clique Sampler's limit.

\subsubsection{QPU Time Scaling}

Pure QPU time scales linearly with number of clusters:
\begin{equation}
t_{\text{QPU}} \approx 40\text{ms} \times n_{\text{clusters}} \times n_{\text{iterations}}
\end{equation}

\subsubsection{Constraint Violations}

Rotation violations are consistently zero after repair, but one-hot violations persist, indicating the penalty $\lambda_{\text{OH}}=3.0$ may be insufficient.

%%%%%%%%%%%%%%%%%%%%%%%%%%%%%%%%%%%%%%%%%%%%%%%%%%%%%%%%%%%%%%%%%%%%%%%%%%%%%%%
\section{Code-to-Equation Verification Summary}
%%%%%%%%%%%%%%%%%%%%%%%%%%%%%%%%%%%%%%%%%%%%%%%%%%%%%%%%%%%%%%%%%%%%%%%%%%%%%%%

\begin{table}[H]
\centering
\caption{Mapping of Variant B equations to implementation}
\begin{adjustbox}{max width=\textwidth}
\begin{tabular}{p{4cm}p{5cm}p{4cm}}
\toprule
\textbf{Variant B Term} & \textbf{Mathematical Form} & \textbf{Code Implementation} \\
\midrule
Base Benefit & $\frac{B_c \cdot L_f}{A_{\text{total}}} \cdot Y_{f,c,t}$ & \texttt{-benefit * area\_frac} \\
Temporal Synergy & $\gamma_{\text{rot}} \cdot R_{c_1,c_2} \cdot \frac{L_f}{A_{\text{total}}}$ & \texttt{-rotation\_gamma * synergy * area\_frac} \\
Spatial Synergy & $\gamma_{\text{spat}} \cdot S_{c_1,c_2} = 0.5\gamma \cdot 0.3R$ & \texttt{-rotation\_gamma * 0.5 * synergy} where \texttt{synergy = R * 0.3} \\
One-Hot Penalty & $2\lambda_{\text{OH}}$ (quadratic term) & \texttt{2 * one\_hot\_penalty} \\
Diversity Bonus & $\lambda_{\text{div}} / T$ & \texttt{diversity\_bonus / n\_periods} \\
Monoculture Penalty & $R_{c,c} = -\beta \cdot 1.5 = -1.2$ & \texttt{negative\_strength * 1.5} \\
\bottomrule
\end{tabular}
\end{adjustbox}
\end{table}

All code implementations correctly match the Variant B formulation from the main project report.

%%%%%%%%%%%%%%%%%%%%%%%%%%%%%%%%%%%%%%%%%%%%%%%%%%%%%%%%%%%%%%%%%%%%%%%%%%%%%%%
\section{Recommendations}
%%%%%%%%%%%%%%%%%%%%%%%%%%%%%%%%%%%%%%%%%%%%%%%%%%%%%%%%%%%%%%%%%%%%%%%%%%%%%%%

\subsection{Cluster Size Guidelines}

\begin{enumerate}
    \item For \textbf{6-family mode}: Use $k \leq 9$ farms/cluster (162 variables max)
    \item For \textbf{27-food hybrid}: Use $k = 2$ farms/cluster (162 variables)
    \item Add safety margin: target 160 variables to avoid edge cases
\end{enumerate}

\subsection{Mode Selection}

\begin{table}[H]
\centering
\caption{Recommended QPU mode by scenario type}
\begin{tabular}{lll}
\toprule
\textbf{Scenario Type} & \textbf{Recommended Mode} & \textbf{Rationale} \\
\midrule
Small 6-food ($\leq$10 farms) & qpu-native-6-family & Direct embedding \\
Medium 6-food (11--200 farms) & qpu-hierarchical-aggregated & Decomposition \\
Any 27-food & qpu-hierarchical-aggregated & 27$\to$6 aggregation \\
Very large (500+ farms) & qpu-hierarchical-aggregated & Scales via clustering \\
\bottomrule
\end{tabular}
\end{table}

\end{document}
